\documentclass[12pt, a4paper]{article}

%<*preamble>
% Math symbols
\usepackage{amsmath, amsthm, amsfonts, amssymb}
\usepackage{accents}
\usepackage{esvect}
\usepackage{mathrsfs}
\usepackage{mathtools}
\mathtoolsset{showonlyrefs}
\usepackage{cmll}
\usepackage{stmaryrd}
\usepackage{physics}
\usepackage[normalem]{ulem}
\usepackage{ebproof}
\usepackage{extarrows}

% Page layout
\usepackage{geometry, a4wide, parskip, fancyhdr}

% Font, encoding, russian support
\usepackage[russian]{babel}
\usepackage[sb]{libertine}
\usepackage{xltxtra}

% Listings
\usepackage{listings}
\lstset{basicstyle=\ttfamily,breaklines=true}
\setmonofont[Scale=MatchLowercase]{JetBrains Mono}

% Miscellaneous
\usepackage{array}
\usepackage{booktabs}\renewcommand{\arraystretch}{1.2}
\usepackage{calc}
\usepackage{caption}
\usepackage{subcaption}
\captionsetup{justification=centering,margin=2cm}
\usepackage{catchfilebetweentags}
\usepackage{enumitem}
\usepackage{etoolbox}
\usepackage{float}
\usepackage{lastpage}
\usepackage{minted}
\usepackage{svg}
\usepackage{wrapfig}
\usepackage{xcolor}
\usepackage[makeroom]{cancel}

\newcolumntype{L}{>{$}l<{$}}
    \newcolumntype{C}{>{$}c<{$}}
\newcolumntype{R}{>{$}r<{$}}

% Footnotes
\usepackage[hang]{footmisc}
\setlength{\footnotemargin}{2mm}
\makeatletter
\def\blfootnote{\gdef\@thefnmark{}\@footnotetext}
\makeatother

% References
\usepackage{hyperref}
\hypersetup{
    colorlinks,
    linkcolor={blue!80!black},
    citecolor={blue!80!black},
    urlcolor={blue!80!black},
}

% tikz
\usepackage{tikz}
\usepackage{tikz-cd}
\usetikzlibrary{arrows.meta}
\usetikzlibrary{decorations.pathmorphing}
\usetikzlibrary{calc}
\usetikzlibrary{patterns}
\usepackage{pgfplots}
\pgfplotsset{width=10cm,compat=1.9}
\newcommand\irregularcircle[2]{% radius, irregularity
    \pgfextra {\pgfmathsetmacro\len{(#1)+rand*(#2)}}
    +(0:\len pt)
    \foreach \a in {10,20,...,350}{
            \pgfextra {\pgfmathsetmacro\len{(#1)+rand*(#2)}}
            -- +(\a:\len pt)
        } -- cycle
}

\providetoggle{useproofs}
\settoggle{useproofs}{false}

\pagestyle{fancy}
\lhead{Лабораторная работа №2}
\lfoot{Михайлов Максим}
\rfoot{M3337}
\cfoot{}
\rhead{стр. \thepage\ из \pageref*{LastPage}}

\newcommand{\R}{\mathbb{R}}
\newcommand{\Q}{\mathbb{Q}}
\newcommand{\Z}{\mathbb{Z}}
\newcommand{\B}{\mathbb{B}}
\newcommand{\N}{\mathbb{N}}
\renewcommand{\Re}{\mathfrak{R}}
\renewcommand{\Im}{\mathfrak{I}}

\newcommand{\const}{\text{const}}
\newcommand{\cond}{\text{cond}}

\newcommand{\teormin}{\textcolor{red}{!}\ }

\DeclareMathOperator*{\xor}{\oplus}
\DeclareMathOperator*{\equ}{\sim}
\DeclareMathOperator{\sign}{\text{sign}}
\DeclareMathOperator{\Sym}{\text{Sym}}
\DeclareMathOperator{\Asym}{\text{Asym}}

\DeclarePairedDelimiter{\ceil}{\lceil}{\rceil}

% godel
\newbox\gnBoxA
\newdimen\gnCornerHgt
\setbox\gnBoxA=\hbox{$\ulcorner$}
\global\gnCornerHgt=\ht\gnBoxA
\newdimen\gnArgHgt
\def\godel #1{%
    \setbox\gnBoxA=\hbox{$#1$}%
    \gnArgHgt=\ht\gnBoxA%
    \ifnum     \gnArgHgt<\gnCornerHgt \gnArgHgt=0pt%
    \else \advance \gnArgHgt by -\gnCornerHgt%
    \fi \raise\gnArgHgt\hbox{$\ulcorner$} \box\gnBoxA %
    \raise\gnArgHgt\hbox{$\urcorner$}}

% \theoremstyle{plain}

\theoremstyle{definition}
\newtheorem{theorem}{Теорема}
\newtheorem*{definition}{Определение}
\newtheorem{axiom}{Аксиома}
\newtheorem*{axiom*}{Аксиома}
\newtheorem{lemma}{Лемма}
\newenvironment{solution}[1][Решение.]{\begin{proof}[#1]}{\end{proof}}

\theoremstyle{remark}
\newtheorem*{remark}{Примечание}
\newtheorem*{exercise}{Упражнение}
\newtheorem{corollary}{Следствие}[theorem]
\newtheorem*{statement}{Утверждение}
\newtheorem*{corollary*}{Следствие}
\newtheorem*{example}{Пример}
\newtheorem{observation}{Наблюдение}
\newtheorem*{prop}{Свойства}
\newtheorem*{obozn}{Обозначение}

% subtheorem
\makeatletter
\newenvironment{subtheorem}[1]{%
    \def\subtheoremcounter{#1}%
    \refstepcounter{#1}%
    \protected@edef\theparentnumber{\csname the#1\endcsname}%
    \setcounter{parentnumber}{\value{#1}}%
    \setcounter{#1}{0}%
    \expandafter\def\csname the#1\endcsname{\theparentnumber.\Alph{#1}}%
    \ignorespaces
}{%
    \setcounter{\subtheoremcounter}{\value{parentnumber}}%
    \ignorespacesafterend
}
\makeatother
\newcounter{parentnumber}

\newtheorem{manualtheoreminner}{Теорема}
\newenvironment{manualtheorem}[1]{%
    \renewcommand\themanualtheoreminner{#1}%
    \manualtheoreminner
}{\endmanualtheoreminner}

\newcommand{\dbltilde}[1]{\accentset{\approx}{#1}}
\newcommand{\intt}{\int\!}

% magical thing that fixes paragraphs
\makeatletter
\patchcmd{\CatchFBT@Fin@l}{\endlinechar\m@ne}{}
{}{\typeout{Unsuccessful patch!}}
\makeatother

\newcommand{\get}[2]{
    \ExecuteMetaData[#1]{#2}
}

\newcommand{\getproof}[2]{
    \iftoggle{useproofs}{\ExecuteMetaData[#1]{#2proof}}{}
}

\newcommand{\getwithproof}[2]{
    \get{#1}{#2}
    \getproof{#1}{#2}
}

\newcommand{\import}[3]{
    \subsection{#1}
    \getwithproof{#2}{#3}
}

\newcommand{\given}[1]{
    Дано выше. (\ref{#1}, стр. \pageref{#1})
}

\renewcommand{\ker}{\text{Ker }}
\newcommand{\im}{\text{Im }}
\renewcommand{\grad}{\text{grad}}
\newcommand{\rg}{\text{rg}}
\newcommand{\defeq}{\stackrel{\text{def}}{=}}
\newcommand{\defeqfor}[1]{\stackrel{\text{def } #1}{=}}
\newcommand{\itemfix}{\leavevmode\makeatletter\makeatother}
\newcommand{\?}{\textcolor{red}{???}}
\renewcommand{\emptyset}{\varnothing}
\newcommand{\longarrow}[1]{\xRightarrow[#1]{\qquad}}
\DeclareMathOperator*{\esup}{\text{ess sup}}
\newcommand\smallO{
    \mathchoice
    {{\scriptstyle\mathcal{O}}}% \displaystyle
    {{\scriptstyle\mathcal{O}}}% \textstyle
    {{\scriptscriptstyle\mathcal{O}}}% \scriptstyle
    {\scalebox{.6}{$\scriptscriptstyle\mathcal{O}$}}%\scriptscriptstyle
}
\renewcommand{\div}{\text{div}\ }
\newcommand{\rot}{\text{rot}\ }
\newcommand{\cov}{\text{cov}}

\makeatletter
\newcommand{\oplabel}[1]{\refstepcounter{equation}(\theequation\ltx@label{#1})}
\makeatother

\newcommand{\symref}[2]{\stackrel{\oplabel{#1}}{#2}}
\newcommand{\symrefeq}[1]{\symref{#1}{=}}

% xrightrightarrows
\makeatletter
\newcommand*{\relrelbarsep}{.386ex}
\newcommand*{\relrelbar}{%
    \mathrel{%
        \mathpalette\@relrelbar\relrelbarsep
    }%
}
\newcommand*{\@relrelbar}[2]{%
    \raise#2\hbox to 0pt{$\m@th#1\relbar$\hss}%
    \lower#2\hbox{$\m@th#1\relbar$}%
}
\providecommand*{\rightrightarrowsfill@}{%
    \arrowfill@\relrelbar\relrelbar\rightrightarrows
}
\providecommand*{\leftleftarrowsfill@}{%
    \arrowfill@\leftleftarrows\relrelbar\relrelbar
}
\providecommand*{\xrightrightarrows}[2][]{%
    \ext@arrow 0359\rightrightarrowsfill@{#1}{#2}%
}
\providecommand*{\xleftleftarrows}[2][]{%
    \ext@arrow 3095\leftleftarrowsfill@{#1}{#2}%
}

\allowdisplaybreaks

\newcommand{\unfinished}{\textcolor{red}{Не дописано}}

% Reproducible pdf builds 
\special{pdf:trailerid [
<00112233445566778899aabbccddeeff>
<00112233445566778899aabbccddeeff>
]}
%</preamble>


\lhead{АиСД, задача 11.2}
\lfoot{Михайлов Максим}
\cfoot{}
\rfoot{M3237}

\begin{document}

\section{Условие}

Научитесь генерировать большие числа Кармайкла с тремя большими простыми делителями.

\section{Решение}

\subsection{Обозначения}

\begin{itemize}
    \item \(\mathcal{C}\) --- множество всех чисел Кармайкла.
    \item \(\mathcal{P}\) --- множество всех простых чисел.
\end{itemize}

\subsection{Теория}

\begin{wrapfigure}{R}{0.5\textwidth}
    \begin{center}
        \includegraphics[width=0.5\textwidth]{images/motivation.jpg}
        \caption{Кот, который мотивировал меня заниматься этой страшной задачей.}
    \end{center}
\end{wrapfigure}

\begin{theorem}\itemfix
    \begin{itemize}
        \item \(\forall i \ \ p_i \in \mathcal{P}\)
        \item \(n = \prod_i p_i\)
        \item \(\forall i \ \ p_i - 1 \equiv 0 \mod n - 1\)
    \end{itemize}
    Тогда \(n \in \mathcal{C}\). В обратную сторону также верно.
\end{theorem}
\begin{proof}
    \(\sphericalangle a : gcd(a, n) = 1\).

    По теореме Ферма:
    \begin{align*}
        a^{p_i - 1} & \equiv 1 \mod p_j   \\
        p_i - 1     & \equiv 0 \mod n - 1 \\
        a^{n - 1}   & \equiv 1 \mod p_j   \\
        a^{n - 1}   & \equiv 1 \mod n
    \end{align*}
\end{proof}

\begin{theorem}\itemfix
    \begin{itemize}
        \item \(p \in \mathcal{P}\)
        \item \(n = pu\)
    \end{itemize}

    Тогда \(p - 1 \equiv 0 \mod n - 1 \Leftrightarrow p - 1 \equiv 0 \mod u - 1\).
\end{theorem}
\begin{proof}
    \[(n - 1) - (u - 1) = n - u = pu - u = (p - 1)u\]
    Если \(p - 1\) делит ровно одно из \(\{n - 1, u - 1\}\), то левая часть не делится на \(p - 1\), а правая делится --- противоречие.
\end{proof}

\subsection{Первый метод}

Будем искать \(n \in \mathcal{C} : n = p \cdot q \cdot r, p < q < r\). По теореме 1:
\[\begin{cases}
        p - 1 \equiv 0 \mod n - 1 \\
        q - 1 \equiv 0 \mod n - 1 \\
        r - 1 \equiv 0 \mod n - 1 \\
    \end{cases}\]
По теореме 2:
\begin{equation}
    \begin{cases}
        p - 1 \equiv 0 \mod qr - 1 \\
        q - 1 \equiv 0 \mod pr - 1 \\
        r - 1 \equiv 0 \mod pq - 1
    \end{cases}
    \label{база}
\end{equation}

Если зафиксировать \(p, q\), то можно перебрать все \(d \equiv 0 \mod pq - 1\), такие что \(q < d\). Для каждого \(d\) нужно проверить, что \(d + 1 \in \mathcal{P}\). Если это так, третье условие выполнено.

\subsection{Второй метод}

Для \(p = 6N + 1, q = 12N + 1, r = 18N + 1\) условие \eqref{база} выполнено:
\[\begin{cases}
        6N \equiv 0 \mod 216 N^2 + 30N   \\
        12N \equiv 0 \mod 108 N^2 + 24 N \\
        18N \equiv 0 \mod 72 N^2 + 18 N
    \end{cases} \]
и нужно только проверять на простоту. Небольшой трюк: если для \(\tilde{N}\) \(p\) не простое, то \(2\tilde{N}\) и \(3 \tilde{N}\) рассматривать не нужно.

\section{Анализ}

\subsection{Первый метод}

Генерировать все \(p, q < n : p,q\in \mathcal{P}\) занимает \(\mathcal{O}(n / \log \log n)\) и асимптотически мы найдём \(\frac{n^2}{\log^2 n}\) пар, т.е.
\[\frac{\log^2 n}{n^2} \cdot \frac{n}{\log \log n} = \frac{\log^2 n}{n \log \log n}\]
итераций на пару. Я не придумал, как оценить число делителей \(pq - 1\), поэтому дальше не анализируется.

\subsection{Второй метод}

Проверить на простоту детерминированно занимает \(\mathcal{O}(\log^6 n)\) с помощью AKS. Т.к. вероятность случайно выбрать простое число есть \(\Theta\left(\frac{1}{\log n}\right)\), то вероятность выбрать 3 простых числа --- \(\Theta\left(\frac{1}{\log^3 n}\right)\). Таким образом, получается \(3 \cdot \mathcal{O}(\log^6 n) \cdot \Theta(\log^3 n) = \mathcal{O}(\log^9 n)\).

Ещё можно использовать модификацию AKS за \(\mathcal{O}(\log^3 n)\), которая работает, если гипотеза Агравала верна, и если мы выдадим неверный ответ, то заявляем всем, что мы её опровергли :).

% \subsection{Третий метод}

% Запишем \eqref{база} в другом виде:

% \[
%     \begin{cases}
%         a(p - 1) = qr - 1 \\
%         b(q - 1) = pr - 1 \\
%         c(r - 1) = pq - 1
%     \end{cases}
% \]

% Пусть \(\delta = bc - q^2\). При фиксированных \(p, с, \delta\) однозначно задается \(q\):

% \begin{theorem}
%     \[q - 1 = \frac{(p - 1)(p + c)}{\delta}\]
% \end{theorem}
% \begin{proof}
%     \begin{align*}
%         q - 1             & = \frac{pr - 1}{b}                        \\
%                           & = \frac{pr - 1 + p - p}{b}                \\
%                           & = \frac{p(r - 1) + (p - 1)}{b}            \\
%                           & = \frac{\frac{p}{c}(pq - 1) + (p - 1)}{b} \\
%                           & = \frac{\frac{p}{c}(pq - 1) + (p - 1)}{b} \\
%                           & = \frac{p(pq - 1) + c(p - 1)}{bc}         \\
%         bc(q - 1)         & = p(pq - 1) + c(p - 1)                    \\
%                           & = p^2q - p + c(p - 1)                     \\
%                           & = p^2q - p - p^2 + p^2 + c(p - 1)         \\
%                           & = p^2(q - 1) - p + p^2 + c(p - 1)         \\
%         (bc - p^2)(q - 1) & = p(p - 1) + c(p - 1)                     \\
%         q - 1             & = \frac{(p + c)(p - 1)}{bc - p^2}         \\
%     \end{align*}
% \end{proof}

% И следовательно однозначно задается \(r\).

% Казалось бы, 

\end{document}
